% Figure 3: where the money goes. Token classes are what Eq. (9) prices
% differently, so the figure is drawn in those classes rather than in a single
% token total. Coordinates from figures/data/token_classes.dat.

\begin{figure}[t]
\centering
\begin{tikzpicture}[font=\small]
\begin{axis}[
  width=0.92\linewidth, height=5.4cm,
  ybar, bar width=7pt,
  ymode=log, log origin=infty,
  ymin=1, ymax=2200,
  ytick={1,10,100,1000},
  yticklabels={1,10,100,1000},
  ylabel={tokens (millions, log scale)},
  xtick={0,1,2,3},
  xticklabels={control,$k{=}0$,$k{=}2$,$k{=}6$},
  enlarge x limits=0.16,
  grid=major, grid style={black!12},
  legend style={at={(0.98,0.97)}, anchor=north east, font=\scriptsize,
                draw=black!25, fill=white, fill opacity=0.92, text opacity=1},
  legend cell align=left,
]
\addplot[draw=black!60, fill=green!25] table[x=x, y=cached]
  {figures/data/token_classes.dat};
\addlegendentry{cached input (cheapest, $\gamma$)}
\addplot[draw=black!60, fill=blue!30] table[x=x, y=uncached]
  {figures/data/token_classes.dat};
\addlegendentry{uncached input ($\alpha$)}
\addplot[draw=black!60, fill=red!45] table[x=x, y=output]
  {figures/data/token_classes.dat};
\addlegendentry{output (most expensive, $\beta$)}
\end{axis}
\end{tikzpicture}
\caption{Token consumption by billing class across the series, the
decomposition Eq.~\eqref{eq:cost} prices. Bars are grouped, not stacked, and
the axis is logarithmic because the classes differ by two orders of magnitude. The control consumes by far the most
tokens in total, but \TotCacheBase{}M of them are cached reads billed at the
cheapest rate: an append-only prompt is almost entirely served from the prefix
cache. The state arms consume far fewer tokens and pay more for them, because
a state block rewritten each turn forfeits that cache (down to
\TotCacheTailSix{}M) and because maintaining the state moves tokens into the
output class. This is why a token-count comparison and a cost comparison
disagree.}
\label{fig:tokenclasses}
\end{figure}
